\clearpage\section{源程序全文}

以下程序由冻结的附录底稿逐段转换排版，未在论文阶段重新计算或改写算法。

\Needspace{10\baselineskip}\subsection{程序 1：formal\_compute.py}

{\raggedright\noindent 原文件：\path{04-compute/src/formal_compute.py}\par}

本文件在代码实现、调试或图表制作过程中得到 Codex 辅助；参数、模型路线、结果解释及最终核验由参赛队负责。

\begingroup\setstretch{1.0}

\VerbatimInput[fontsize=\fontsize{7.5}{9}\selectfont,breaklines=true,breakanywhere=true,numbers=left,numbersep=5pt]{code/program-01.py}

\endgroup

\Needspace{10\baselineskip}\subsection{程序 2：finalize\_compute.py}

{\raggedright\noindent 原文件：\path{04-compute/src/finalize_compute.py}\par}

本文件在代码实现、调试或图表制作过程中得到 Codex 辅助；参数、模型路线、结果解释及最终核验由参赛队负责。

\begingroup\setstretch{1.0}

\VerbatimInput[fontsize=\fontsize{7.5}{9}\selectfont,breaklines=true,breakanywhere=true,numbers=left,numbersep=5pt]{code/program-02.py}

\endgroup

\Needspace{10\baselineskip}\subsection{程序 3：q3\_dense\_sensitivity.py}

{\raggedright\noindent 原文件：\path{04-compute/src/q3_dense_sensitivity.py}\par}

本文件在代码实现、调试或图表制作过程中得到 Codex 辅助；参数、模型路线、结果解释及最终核验由参赛队负责。

\begingroup\setstretch{1.0}

\VerbatimInput[fontsize=\fontsize{7.5}{9}\selectfont,breaklines=true,breakanywhere=true,numbers=left,numbersep=5pt]{code/program-03.py}

\endgroup

\Needspace{10\baselineskip}\subsection{程序 4：repair\_probe.py}

{\raggedright\noindent 原文件：\path{04-compute/src/repair_probe.py}\par}

本文件在代码实现、调试或图表制作过程中得到 Codex 辅助；参数、模型路线、结果解释及最终核验由参赛队负责。

\begingroup\setstretch{1.0}

\VerbatimInput[fontsize=\fontsize{7.5}{9}\selectfont,breaklines=true,breakanywhere=true,numbers=left,numbersep=5pt]{code/program-04.py}

\endgroup

\Needspace{10\baselineskip}\subsection{程序 5：run\_q2\_grid\_dt5.py}

{\raggedright\noindent 原文件：\path{04-compute/revision-grid-conv-20260913/run_q2_grid_dt5.py}\par}

本文件在代码实现、调试或图表制作过程中得到 Codex 辅助；参数、模型路线、结果解释及最终核验由参赛队负责。

\begingroup\setstretch{1.0}

\VerbatimInput[fontsize=\fontsize{7.5}{9}\selectfont,breaklines=true,breakanywhere=true,numbers=left,numbersep=5pt]{code/program-05.py}

\endgroup

\Needspace{10\baselineskip}\subsection{程序 6：run\_q2\_model\_trajectories.py}

{\raggedright\noindent 原文件：\path{04-compute/revision-grid-conv-20260913/run_q2_model_trajectories.py}\par}

本文件在代码实现、调试或图表制作过程中得到 Codex 辅助；参数、模型路线、结果解释及最终核验由参赛队负责。

\begingroup\setstretch{1.0}

\VerbatimInput[fontsize=\fontsize{7.5}{9}\selectfont,breaklines=true,breakanywhere=true,numbers=left,numbersep=5pt]{code/program-06.py}

\endgroup

\Needspace{10\baselineskip}\subsection{程序 7：run\_space\_extra.py}

{\raggedright\noindent 原文件：\path{04-compute/revision-grid-conv-20260913/run_space_extra.py}\par}

本文件在代码实现、调试或图表制作过程中得到 Codex 辅助；参数、模型路线、结果解释及最终核验由参赛队负责。

\begingroup\setstretch{1.0}

\VerbatimInput[fontsize=\fontsize{7.5}{9}\selectfont,breaklines=true,breakanywhere=true,numbers=left,numbersep=5pt]{code/program-07.py}

\endgroup

\Needspace{10\baselineskip}\subsection{程序 8：redraw\_fig6.py}

{\raggedright\noindent 原文件：\path{06-figure/scripts/redraw_fig6.py}\par}

本文件在代码实现、调试或图表制作过程中得到 Codex 辅助；参数、模型路线、结果解释及最终核验由参赛队负责。

\begingroup\setstretch{1.0}

\VerbatimInput[fontsize=\fontsize{7.5}{9}\selectfont,breaklines=true,breakanywhere=true,numbers=left,numbersep=5pt]{code/program-08.py}

\endgroup

\Needspace{10\baselineskip}\subsection{程序 9：redraw\_fig7.py}

{\raggedright\noindent 原文件：\path{06-figure/scripts/redraw_fig7.py}\par}

本文件在代码实现、调试或图表制作过程中得到 Codex 辅助；参数、模型路线、结果解释及最终核验由参赛队负责。

\begingroup\setstretch{1.0}

\VerbatimInput[fontsize=\fontsize{7.5}{9}\selectfont,breaklines=true,breakanywhere=true,numbers=left,numbersep=5pt]{code/program-09.py}

\endgroup

\Needspace{10\baselineskip}\subsection{程序 10：redraw\_fig9\_onward.py}

{\raggedright\noindent 原文件：\path{06-figure/scripts/redraw_fig9_onward.py}\par}

本文件在代码实现、调试或图表制作过程中得到 Codex 辅助；参数、模型路线、结果解释及最终核验由参赛队负责。

\begingroup\setstretch{1.0}

\VerbatimInput[fontsize=\fontsize{7.5}{9}\selectfont,breaklines=true,breakanywhere=true,numbers=left,numbersep=5pt]{code/program-10.py}

\endgroup

\Needspace{10\baselineskip}\subsection{程序 11：redraw\_q1\_end\_effect\_time.py}

{\raggedright\noindent 原文件：\path{06-figure/scripts/redraw_q1_end_effect_time.py}\par}

本文件在代码实现、调试或图表制作过程中得到 Codex 辅助；参数、模型路线、结果解释及最终核验由参赛队负责。

\begingroup\setstretch{1.0}

\VerbatimInput[fontsize=\fontsize{7.5}{9}\selectfont,breaklines=true,breakanywhere=true,numbers=left,numbersep=5pt]{code/program-11.py}

\endgroup

\Needspace{10\baselineskip}\subsection{程序 12：render\_q2\_profiles.py}

{\raggedright\noindent 原文件：\path{06-figure/scripts/render_q2_profiles.py}\par}

本文件在代码实现、调试或图表制作过程中得到 Codex 辅助；参数、模型路线、结果解释及最终核验由参赛队负责。

\begingroup\setstretch{1.0}

\VerbatimInput[fontsize=\fontsize{7.5}{9}\selectfont,breaklines=true,breakanywhere=true,numbers=left,numbersep=5pt]{code/program-12.py}

\endgroup

\Needspace{10\baselineskip}\subsection{程序 13：render\_q2\_profiles\_3d.py}

{\raggedright\noindent 原文件：\path{06-figure/scripts/render_q2_profiles_3d.py}\par}

本文件在代码实现、调试或图表制作过程中得到 Codex 辅助；参数、模型路线、结果解释及最终核验由参赛队负责。

\begingroup\setstretch{1.0}

\VerbatimInput[fontsize=\fontsize{7.5}{9}\selectfont,breaklines=true,breakanywhere=true,numbers=left,numbersep=5pt]{code/program-13.py}

\endgroup

\Needspace{10\baselineskip}\subsection{程序 14：revise\_q2\_figures\_20260913.py}

{\raggedright\noindent 原文件：\path{06-figure/scripts/revise_q2_figures_20260913.py}\par}

本文件在代码实现、调试或图表制作过程中得到 Codex 辅助；参数、模型路线、结果解释及最终核验由参赛队负责。

\begingroup\setstretch{1.0}

\VerbatimInput[fontsize=\fontsize{7.5}{9}\selectfont,breaklines=true,breakanywhere=true,numbers=left,numbersep=5pt]{code/program-14.py}

\endgroup

\Needspace{10\baselineskip}\subsection{程序 15：update\_figure\_plan\_20260913.py}

{\raggedright\noindent 原文件：\path{06-figure/scripts/update_figure_plan_20260913.py}\par}

本文件在代码实现、调试或图表制作过程中得到 Codex 辅助；参数、模型路线、结果解释及最终核验由参赛队负责。

\begingroup\setstretch{1.0}

\VerbatimInput[fontsize=\fontsize{7.5}{9}\selectfont,breaklines=true,breakanywhere=true,numbers=left,numbersep=5pt]{code/program-15.py}

\endgroup

\Needspace{10\baselineskip}\subsection{程序 16：verify\_current\_pipeline.py}

{\raggedright\noindent 原文件：\path{06-figure/scripts/verify_current_pipeline.py}\par}

本文件在代码实现、调试或图表制作过程中得到 Codex 辅助；参数、模型路线、结果解释及最终核验由参赛队负责。

\begingroup\setstretch{1.0}

\VerbatimInput[fontsize=\fontsize{7.5}{9}\selectfont,breaklines=true,breakanywhere=true,numbers=left,numbersep=5pt]{code/program-16.py}

\endgroup

\Needspace{10\baselineskip}\subsection{程序 17：make\_official\_contact\_sheet.py}

{\raggedright\noindent 原文件：\path{06-figure/scripts/make_official_contact_sheet.py}\par}

本文件在代码实现、调试或图表制作过程中得到 Codex 辅助；参数、模型路线、结果解释及最终核验由参赛队负责。

\begingroup\setstretch{1.0}

\VerbatimInput[fontsize=\fontsize{7.5}{9}\selectfont,breaklines=true,breakanywhere=true,numbers=left,numbersep=5pt]{code/program-17.py}

\endgroup
